% Einheit 4 von 5 – Flächenbedingungen (bank/flaecheninhalt-durch-integration/e4.jsonl, zusammenbau v0.1)
%% TODO Zeitmarke Einheit 4: Marken-Zeile nicht lesbar
%% TODO Zweigzeile Teil 1: was hier gelernt wird (Satz in Schülersprache aus dem Einheitstitel) – Einheit 4
\einheitenkopf[e4]{Einheit 4 von 5 · Flächenbedingungen}
\zweigzeile{Abitur GK}
\setcounter{aufgabe}{14}

%% TODO Titel: Ich-kann-Satz fehlt in der Bank (Platzhalter: Kettenname) – Nr. 15
%% TODO Anweisung: ein Satz über der ersten Teilaufgabe fehlt in der Bank – Nr. 15
\begin{aufgabe}{Flächenbedingungen}
%% TODO \rechenplatz{n}: Schrittzahl des Grundfalls fehlt in der Bank (nur bei fester Schreibform, 2.3 b)
\begin{teile}
\teil Gesucht ist der Inhalt der Fläche zwischen dem Graphen von $f(x) = x^2$ und der Geraden $y = 4$ – vorwärts oder rückwärts? \kreuz{vorwärts (Inhalt gesucht)}\\ \kreuz{rückwärts (Inhalt gegeben)}
\teil Für $m > 0$ schließt die Gerade $y = m \cdot x$ mit dem Graphen von $f(x) = 2x^2$ eine Fläche vom Inhalt $9$ ein – $m$?
\teil Eine Fläche liegt über der x-Achse rechts der y-Achse: bis $x = 2$ unter der Geraden $y = 3$, danach unter der Strecke von $(2|3)$ zum Punkt $(b|0)$ der x-Achse. Sie hat den Inhalt $12$ – $b$?
\teil Die Fläche unter dem Graphen von $f(x) = 3x^2$ von $0$ bis $2$ soll von einer Parallelen zur y-Achse $x = k$ halbiert werden – $k$? Runde auf zwei Nachkommastellen.
\teil Für $a > 0$ schließt der Graph von $f_a(x) = -a x^2 + 4$ mit der x-Achse eine Fläche vom Inhalt $16$ ein – $a$?
\teil $f(x) = e^{x}$. Begründe ohne Rechnung, dass es genau ein $b > 0$ gibt, für das die Fläche unter dem Graphen von $0$ bis $b$ den Inhalt $5$ hat.
\teil Der Graph von $f$ in der Abbildung ist punktsymmetrisch zu $(0|3)$. Die Gleichung Integral von $k$ bis $k + 2$ über $f(x)$ $= 6$ hat genau eine Lösung mit $-1{,}5 \le k \le 0$. Gib sie ohne Rechnung an und begründe mit Eintragungen in der Abbildung. \hfill (Abitur 2025 GK)

\begin{ksys}[xmin=-3,xmax=3,ymin=-1,ymax=6,ablesen]\funktion{0.5*\x^3-2*\x+3}{f}\end{ksys}
\end{teile}
\end{aufgabe}

%% TODO Titel: Ich-kann-Satz fehlt in der Bank (Platzhalter: Kettenname) – Nr. 16
%% TODO Anweisung: ein Satz über der ersten Teilaufgabe fehlt in der Bank – Nr. 16
\begin{aufgabe}{Flächenbedingungen – Fehler finden, Begründen, Anwendung}
\begin{teile}
\teil Die Fläche unter dem Graphen von $f(x) = 0{,}75x^2$ von $0$ bis $4$ soll durch $x = k$ halbiert werden – wo ist der Fehler? \rechnung{f(k) &= \tfrac{1}{2} f(4) = 6 \\ 0{,}75k^2 &= 6 \\ k &\approx 2{,}83}
\teil Warum nimmt ein stetiger, streng monoton wachsender Flächenterm jeden Wert zwischen Anfangs- und Endwert genau einmal an?
\teil Ein Grundstück liegt zwischen Straße (x-Achse), Mauer (y-Achse), der Geraden $x = 4$ und einem Bach entlang $f(x) = 0{,}25x^2 + 2$; eine Einheit entspricht $10$ m. Zwei Erben teilen es durch einen Zaun parallel zur Mauer in gleich große Teile – wie weit von der Mauer steht der Zaun?
\end{teile}
\end{aufgabe}
